\documentclass[11pt,letterpaper]{letter}
\usepackage[margin=1.2in]{geometry}
\usepackage{microtype}
\usepackage{xcolor}
\usepackage[colorlinks=true, urlcolor=blue!70!black,
            linkcolor=blue!70!black, citecolor=blue!70!black]{hyperref}
\usepackage{bookmark}

\signature{Aryan Hemendra Shah\\
  University of Miami\\
  \href{mailto:ahs222@miami.edu}{ahs222@miami.edu}\\
  ORCID: 0009-0008-0518-428X}

\address{Aryan Hemendra Shah\\
  University of Miami\\
  Miami, FL 33146\\
  \href{mailto:ahs222@miami.edu}{ahs222@miami.edu}}

\date{\today}

\begin{document}

\begin{letter}{The Editors\\
  \textit{Life Sciences in Space Research}\\
  Elsevier}

\opening{Dear Editors,}

I am submitting for your consideration the manuscript
\textbf{``An open-source GCR dosimetry pipeline for Mars mission risk
assessment: organ-specific REID, SEP acute hazard, and dual-use proton
therapy RBE benchmarking''} for publication in \textit{Life Sciences in
Space Research}.

\medskip
\textbf{What the paper does.}
The manuscript presents an open-source Python package
(\textsc{gcr-dosimetry-pipeline}) that integrates GCR spectrum
generation, CSDA slab transport, organ-specific cancer risk (REID),
Solar Energetic Particle acute dosimetry, and global uncertainty
quantification via Latin Hypercube Sampling into a single reproducible
workflow with 82 automated tests and a public repository.

\medskip
\textbf{What is new.}
I am transparent that the individual components — force-field
modulation, CSDA transport, REID estimation — are established methods.
The contributions are threefold.

First, no currently available open tool combines organ-specific REID,
SEP acute risk, and global uncertainty propagation in a single workflow
reproducible from a standard \texttt{pip install}.

Second, the per-species LET decomposition shows that carbon ions account
for approximately 59\% of unshielded GCR absorbed dose at a field LET
of 18.7~keV/$\mu$m — the same LET range that governs Bragg-peak RBE
uncertainty in clinical proton therapy.
The pipeline computes Wedenberg and McNamara RBE for both GCR and
clinical proton beam geometries from one codebase, enabling direct
comparison of quality factors across both fields.
To my knowledge this connection has not been demonstrated quantitatively
in an open tool.

Third, the Latin Hypercube Sampling ensemble shows that DDREF and excess
risk scaling together account for more than 70\% of REID output variance,
quantifying why physics improvements alone cannot close mission risk
uncertainty. This is a directly actionable finding for research
prioritization.

\medskip
\textbf{Scope and limitations.}
The transport model uses CSDA with exponential nuclear attenuation and
does not capture heavy-ion fragmentation secondaries.
This limitation is stated explicitly and its domain of validity is
discussed.
The per-species calibration to MSL/RAD is described as calibration, not
independent prediction, and the validation evidence is presented
accordingly.
The REID framework follows Cucinotta (2013); the 2023 NASA update is
noted as a limitation for future work.

\medskip
\textbf{Relevance to \textit{Life Sciences in Space Research}.}
This manuscript is squarely within the journal's scope: human
spaceflight radiation risk, Mars mission planning, and space radiation
biology.
The dual-use LET framing (space radiation $\leftrightarrow$ proton
therapy) additionally connects to an emerging readership interested in
high-LET radiobiology across both domains.

\medskip
\textbf{Data and code availability.}
All source code, input data, validation scripts, and TOPAS benchmark
geometries are openly available on GitHub at the repository
\href{https://github.com/aryanhshah8/gcr-dosimetry-pipeline}{aryanhshah8/gcr-dosimetry-pipeline},
and no proprietary datasets are required to reproduce any results.

\medskip
I confirm that this manuscript has not been submitted elsewhere, that I
am the sole author, and that I have no conflicts of interest to declare.
I look forward to the reviewers' comments.

\closing{Sincerely,}

\end{letter}
\end{document}
